% Biên dịch bằng XeLaTeX:  xelatex bao_cao_ai_training.tex  (chạy trong docs/report để thấy hình)
\documentclass[12pt,a4paper]{article}
\usepackage{fontspec}
\setmainfont{TeX Gyre Termes}
\usepackage[a4paper,margin=2.5cm]{geometry}
\usepackage{booktabs,array,longtable,graphicx,float,enumitem,amsmath,xcolor}
\usepackage[hidelinks]{hyperref}
\setlist{nosep,leftmargin=1.5em}
\renewcommand{\arraystretch}{1.15}
\newcommand{\fig}[2]{\IfFileExists{#1}{\begin{figure}[H]\centering\includegraphics[width=0.85\linewidth]{#1}\caption{#2}\end{figure}}{}}

\title{\textbf{BÁO CÁO CHI TIẾT DỰ ÁN}\\[4pt]
\large XÂY DỰNG HỆ THỐNG HỖ TRỢ GỢI Ý ĐƯỜNG ĐI CỦA CƠN BÃO VIỆT NAM\\[2pt]
\normalsize Trọng tâm: huấn luyện và đánh giá mô hình trí tuệ nhân tạo}
\author{Typhoon VN Forecast System}
\date{06/11/2026}

\begin{document}
\maketitle

\section{Thông tin chung}
\begin{tabular}{@{}ll@{}}
Tên dự án & Typhoon VN Forecast System \\
Phiên bản khai báo & 0.1.0 \\
Ngôn ngữ và nền tảng & Python 3.11, PyTorch, scikit-learn, Pandas, NumPy, FastAPI \\
Giấy phép & MIT \\
Thời điểm đánh giá & 06/11/2026 \\
\end{tabular}

\medskip
Mục tiêu của dự án là xây dựng hệ thống AI hỗ trợ ra quyết định, gợi ý đường đi
của xoáy thuận nhiệt đới có khả năng ảnh hưởng Việt Nam và Biển Đông. Báo cáo này
tập trung vào phần huấn luyện mô hình: cách phát biểu bài toán, dữ liệu huấn luyện,
các phương pháp AI đã so sánh, quy trình chọn mô hình, kết quả thực nghiệm và
các thí nghiệm tối ưu huấn luyện.

\textbf{Giới hạn an toàn:} Sản phẩm chỉ mang tính tham khảo, phục vụ học tập và
nghiên cứu; không phải bản tin dự báo khí tượng chính thức và không thay thế bản
tin của Trung tâm Dự báo Khí tượng Thủy văn Quốc gia (NCHMF).

\section{Bài toán và phạm vi}
``Gợi ý đường đi của cơn bão'' được mô hình hóa là \textbf{bài toán dự báo chuỗi
thời gian nhiều bước} (multi-horizon regression), không phải bài toán tìm đường:
cơn bão không ``chọn'' đường đi, nên mục tiêu là ước lượng vị trí tâm bão tương lai
cùng độ bất định của nó.

\begin{itemize}
  \item \textbf{Đầu vào:} 4 quan trắc gần nhất cách nhau 6 giờ (24 giờ lịch sử);
        mỗi quan trắc có 47 đặc trưng.
  \item \textbf{Đầu ra:} vị trí tâm bão tại $+6, +12, +24, +48, +72$ giờ.
  \item \textbf{Biểu diễn đầu ra:} \emph{độ dịch chuyển} $(\Delta\varphi, \Delta\lambda)$
        so với vị trí cuối, thay vì tọa độ tuyệt đối (mục~\ref{sec:ablation}).
  \item \textbf{Hàm mất mát:} Smooth~L1 trên độ dịch chuyển (đơn vị độ).
  \item \textbf{Độ đo:} sai số khoảng cách vòng lớn (Haversine, km) giữa vị trí dự
        báo và vị trí thực; độ phủ (coverage) của nón bất định.
\end{itemize}

\noindent\textbf{Trong phạm vi:} chuẩn bị dữ liệu, tạo đặc trưng, huấn luyện và so
sánh các phương pháp AI, kiểm định thống kê, đóng gói mô hình phục vụ API/demo.\\
\textbf{Ngoài phạm vi:} phát hành cảnh báo chính thức, dự báo cường độ/mưa/lũ,
gửi thông báo thực tế.

\section{Kiến trúc tổng thể pipeline huấn luyện}
\begin{center}
\small IBTrACS WP (1980--2024) $\rightarrow$ làm sạch, chuẩn hóa $\rightarrow$
đặc trưng quỹ đạo (+ dòng dẫn đường NCEP/NCAR) $\rightarrow$ cửa sổ chuỗi, chia theo năm
$\rightarrow$ huấn luyện (validation) $\rightarrow$ đánh giá test một lần
$\rightarrow$ model bundle $\rightarrow$ FastAPI / dashboard.
\end{center}
Các module độc lập: \texttt{ingestion} (tải, parse), \texttt{features} (làm sạch,
đặc trưng, scaler), \texttt{datasets} (cửa sổ, chia tập), \texttt{models} và
\texttt{training} (baseline, LSTM, Transformer, checkpoint), \texttt{evaluation}
(sai số track, bootstrap), \texttt{inference} (nạp bundle, k-NN analog, nón bất định).

\section{Dữ liệu và quản trị nguồn gốc}
\begin{itemize}
  \item \textbf{Nguồn:} NOAA IBTrACS v04r01, lưu vực Tây Bắc Thái Bình Dương (WP),
        giai đoạn 1980--2024; tọa độ WGS84, thời gian UTC.
  \item \textbf{Làm sạch:} chuẩn hóa đơn vị (knot $\rightarrow$ m/s), khử trùng lặp,
        chỉ giữ quan trắc synop 6 giờ, giới hạn vùng 0--30°N, 100--140°E.
        Sau làm sạch còn \textbf{24.024 quan trắc, 1.089 cơn bão}.
  \item \textbf{Dữ liệu môi trường:} tái phân tích NCEP/NCAR R1 (lưới 2,5°, miễn phí)
        lấy đúng thời điểm quan trắc để tạo đặc trưng dòng dẫn đường.
  \item \textbf{Truy vết:} mỗi quan sát lưu nguồn, URL và SHA-256 tệp gốc; dữ liệu
        vi phạm WGS84, thiếu múi giờ hoặc giá trị không hữu hạn bị từ chối.
\end{itemize}

\subsection*{Chia tập không rò rỉ}
Chia \textbf{theo năm, giữ nguyên cả cơn bão}; scaler chỉ fit trên train.

\begin{table}[H]\centering
\begin{tabular}{lccc}
\toprule
Tập & Giai đoạn & Số bão & Ghi chú \\
\midrule
Train & 1980--2016 & 895 & $\approx$7.800 cửa sổ huấn luyện \\
Validation & 2017--2020 & 110 & chọn mô hình, siêu tham số \\
Test & 2021--2024 & 84 & chạy một lần, 714 cửa sổ \\
Tập con ảnh hưởng VN & 2021--2024 & 27 & đi vào trong 300 km bờ biển, 278 cửa sổ \\
\bottomrule
\end{tabular}
\caption{Chia tập dữ liệu}
\end{table}

\section{Làm sạch và tạo đặc trưng}
Một \texttt{FeatureBuilder} dùng chung cho huấn luyện và suy luận (tránh
training--serving skew). 47 đặc trưng theo thứ tự cố định:
\begin{itemize}
  \item \textbf{Vị trí và độ trễ:} vĩ độ, kinh độ, chênh lệch và lag 1--3 bước.
  \item \textbf{Động học:} quãng đường bước, hướng (bearing), tốc độ km/h, gia tốc,
        góc rẽ, tốc độ rẽ; timestamp trùng/lùi không sinh tốc độ giả.
  \item \textbf{Thời gian:} sin/cos của giờ, ngày, tháng, ngày trong năm; cờ mùa bão.
  \item \textbf{Địa lý:} khoảng cách, hướng tới bờ biển Việt Nam và các điểm tham
        chiếu (Hoàng Sa, Trường Sa, Đà Nẵng, \dots).
  \item \textbf{Cường độ:} gió (m/s), khí áp (hPa), biến thiên và one-hot cấp bão.
\end{itemize}
\textbf{Đặc trưng dòng dẫn đường (11 biến, thí nghiệm mục~\ref{sec:steering}):}
dòng dẫn đường lớp sâu trong vành 200--800 km (\texttt{dlm\_u}, \texttt{dlm\_v}),
gió 500 hPa, đứt gió 200--850 hPa, độ cao địa thế vị 500 hPa và chênh lệch 5° về
bốn hướng (phản ánh vị trí áp cao cận nhiệt).

\texttt{FeatureScaler} chỉ fit trên train, lưu thứ tự cột, giá trị bù, mean, scale;
từ chối schema thiếu/thừa/đổi thứ tự cột và xử lý cột toàn NaN/Inf.

\section{Mô hình và huấn luyện}
\subsection{Các phương pháp so sánh}
\begin{table}[H]\centering\small
\begin{tabular}{l l p{8cm}}
\toprule
Nhóm & Phương pháp & Ý tưởng \\
\midrule
Mốc so sánh & Persistence & Giữ nguyên vận tốc 6 giờ gần nhất \\
Mốc so sánh & CLIPER-proxy & Persistence có suy giảm (không phải CLIPER5 chính thức) \\
AI cổ điển & k-NN analog & Lập luận dựa trên tình huống: tìm các đoạn bão lịch sử
giống nhất, trung bình có trọng số đường đi tiếp theo, \textbf{nêu tên bão tương tự} để giải thích \\
Học máy & Gradient boosting & HistGradientBoosting hồi quy độ dịch chuyển \\
Học sâu & LSTM & 2 lớp, hidden 128, dropout 0,2; đọc 4 bước, đầu ra 5 mốc \\
Học sâu & Attention-LSTM & LSTM + attention cộng trên chuỗi đầu vào \\
Học sâu & Transformer & Encoder self-attention + mã hóa vị trí \\
\bottomrule
\end{tabular}
\caption{Các phương pháp AI được huấn luyện và so sánh}
\end{table}

\subsection{Cấu hình huấn luyện}
\begin{itemize}
  \item Bộ tối ưu Adam, learning rate $10^{-3}$, weight decay $10^{-5}$,
        lịch cosine, gradient clipping 1,0.
  \item Tối đa 100 epoch, early stopping patience 10, giữ trọng số epoch tốt nhất
        trên validation; hỗ trợ checkpoint và resume.
  \item Mạng nơ-ron chạy 3 seed (0, 1, 2), báo cáo trung bình $\pm$ độ lệch chuẩn;
        persistence, CLIPER và k-NN là tất định.
  \item Siêu tham số k-NN chọn trên validation: $k=40$, trọng số vị trí 1,0, mùa 0,5.
  \item Seq2Seq có trong mã nguồn nhưng không đưa vào so sánh do teacher forcing
        trộn tọa độ chuẩn hóa với mục tiêu vật lý.
\end{itemize}

\subsection{Quy trình chọn mô hình}
Mọi lựa chọn (kiến trúc, đặc trưng, siêu tham số, thành viên ensemble) chỉ dựa trên
\textbf{validation 2017--2020}. Tập test 2021--2024 được chạy \textbf{một lần}
sau khi đã chốt cấu hình. Kiểm định ý nghĩa dùng bootstrap theo \emph{cơn bão}
(các dự báo của cùng một bão tương quan với nhau), 5.000 lần lấy mẫu lại.

\section{Model bundle và tính tái lập}
Mô hình sau huấn luyện được đóng gói gồm \texttt{weights.pt}, \texttt{scaler.json},
\texttt{split-manifest.json} và \texttt{manifest.json} (phiên bản, SHA-256 dataset,
cấu hình mô hình, độ dài đầu vào, horizon, hash schema đặc trưng, bán kính nón,
checksum từng tệp). Khi nạp: allowlist loại mô hình, kiểm tra checksum, kích thước
đặc trưng, horizon và scaler; PyTorch nạp với \texttt{weights\_only=True}.

Tái lập kết quả:
{\small
\begin{verbatim}
python scripts/run_ibtracs_benchmark.py --data-root data \
    --output-dir runs/ibtracs-benchmark
python scripts/run_model_comparison.py --data-root data \
    --output-dir runs/model-comparison --seeds 0 1 2
python scripts/plot_error_analysis.py --run runs/model-comparison \
    --data-root data
python scripts/build_analog_artifact.py --data-root data \
    --output runs/analog/analog-wp.npz
\end{verbatim}}

\section{Phục vụ mô hình (API và demo)}
\begin{itemize}
  \item FastAPI: \texttt{/forecast}, \texttt{/typhoons/active},
        \texttt{/typhoons/\{id\}/track} (dự báo + nón), \texttt{/health},
        \texttt{/ready}, \texttt{/version}; worker lấy dữ liệu NOAA IBTrACS ACTIVE.
  \item Mỗi dự báo kèm đường đi gợi ý của k-NN analog và 5 cơn bão lịch sử giống nhất
        (tên, năm). Thư viện analog chỉ gồm bão 1980--2020, không chứa bão test.
  \item Demo công khai \url{https://baoviet2026.pages.dev}: 20 bão thật 2021--2024
        ảnh hưởng Việt Nam; sai số 72 h trung bình của LSTM trên 20 bão: 408 km.
\end{itemize}

\section{Hạ tầng phát triển}
Khai báo dependency trong \texttt{pyproject.toml} (nhóm \texttt{train}: PyTorch,
Pandas, scikit-learn; \texttt{experiment}: DVC, MLflow). Docker Compose cung cấp
API, worker, PostGIS (migration Alembic, truy vấn bán kính \texttt{ST\_DWithin}) và
MinIO; image Docker phục vụ mô hình đã huấn luyện, chạy với user không phải root.

\section{Kết quả thực nghiệm}
\subsection{Sai số trung bình trên toàn bộ 84 bão test (km)}
\begin{table}[H]\centering
\begin{tabular}{lrrrrr}
\toprule
Phương pháp & 6h & 12h & 24h & 48h & 72h \\
\midrule
Persistence & 31 & 68 & 155 & 390 & 676 \\
CLIPER-proxy & 30 & 65 & 151 & 388 & 677 \\
k-NN analog & 32 & 64 & 133 & 290 & \textbf{462} \\
Gradient boosting & \textbf{29} & \textbf{60} & \textbf{129} & 292 & 476 \\
LSTM & 33 $\pm$ 0 & 65 $\pm$ 1 & 133 $\pm$ 1 & \textbf{289} $\pm$ 1 & 464 $\pm$ 5 \\
Attention-LSTM & 35 $\pm$ 3 & 68 $\pm$ 4 & 138 $\pm$ 5 & 297 $\pm$ 10 & 473 $\pm$ 9 \\
Transformer & 33 $\pm$ 1 & 64 $\pm$ 1 & 133 $\pm$ 3 & 293 $\pm$ 3 & 468 $\pm$ 5 \\
\bottomrule
\end{tabular}
\caption{Sai số track trung bình, toàn tập test 2021--2024}
\end{table}

\subsection{Sai số trên 27 bão ảnh hưởng Việt Nam (km)}
\begin{table}[H]\centering
\begin{tabular}{lrrrrr}
\toprule
Phương pháp & 6h & 12h & 24h & 48h & 72h \\
\midrule
Persistence & 33 & 70 & 159 & 388 & 642 \\
CLIPER-proxy & 32 & 68 & 159 & 414 & 710 \\
k-NN analog & 33 & 65 & 134 & 290 & 476 \\
Gradient boosting & \textbf{30} & 61 & 131 & 295 & 491 \\
LSTM & 32 $\pm$ 0 & 63 $\pm$ 0 & 133 $\pm$ 1 & 292 $\pm$ 6 & 468 $\pm$ 9 \\
Attention-LSTM & 34 $\pm$ 3 & 66 $\pm$ 4 & 138 $\pm$ 4 & 298 $\pm$ 6 & 471 $\pm$ 3 \\
Transformer & 32 $\pm$ 1 & \textbf{62} $\pm$ 1 & \textbf{130} $\pm$ 2 & \textbf{289} $\pm$ 3 & \textbf{463} $\pm$ 10 \\
\bottomrule
\end{tabular}
\caption{Sai số track trung bình, bão ảnh hưởng Việt Nam}
\end{table}

\fig{figures/error_by_horizon.png}{Sai số theo thời hạn dự báo}

\subsection{Kiểm định thống kê (bootstrap theo cơn bão, 72 h, 50 bão)}
\begin{table}[H]\centering\small
\begin{tabular}{lrr}
\toprule
Phương pháp & So với persistence (km, KTC 95\%) & So với LSTM (km, KTC 95\%) \\
\midrule
CLIPER-proxy & $-15$ [$-101$, $+65$] & $+188$ [$+113$, $+256$] \\
k-NN analog & $-201$ [$-268$, $-131$] & $+3$ [$-29$, $+34$] \\
Gradient boosting & $-191$ [$-252$, $-126$] & $+12$ [$-20$, $+48$] \\
LSTM & $-203$ [$-262$, $-145$] & --- \\
Attention-LSTM & $-173$ [$-229$, $-116$] & $+30$ [$+12$, $+50$] \\
Transformer & $-200$ [$-257$, $-142$] & $+3$ [$-16$, $+22$] \\
\bottomrule
\end{tabular}
\caption{Chênh lệch sai số 72 h (âm = tốt hơn mốc tham chiếu)}
\end{table}
Mọi phương pháp học tốt hơn persistence 170--200 km có ý nghĩa thống kê. Giữa
LSTM, Transformer, k-NN và gradient boosting \textbf{không có khác biệt có ý nghĩa};
Attention-LSTM kém hơn LSTM có ý nghĩa.

\subsection{Nón bất định}
Bán kính nón lấy phân vị 80\% sai số trên validation: 52 / 97 / 210 / 455 / 714 km
tại 6 / 12 / 24 / 48 / 72 h. Độ phủ thực tế trên test đạt 82--86\%, sát mục tiêu 80\%.

\subsection{Thí nghiệm loại bỏ: mục tiêu học độ dịch chuyển}\label{sec:ablation}
Cùng kiến trúc LSTM, cùng dữ liệu, seed 0, chỉ đổi mục tiêu học:
\begin{table}[H]\centering
\begin{tabular}{lrrrrr}
\toprule
Mục tiêu của LSTM & 6h & 12h & 24h & 48h & 72h \\
\midrule
Tọa độ tuyệt đối & 115 & 121 & 164 & 305 & 473 \\
Độ dịch chuyển (đang dùng) & \textbf{33} & \textbf{65} & \textbf{133} & \textbf{289} & \textbf{464} \\
\bottomrule
\end{tabular}
\end{table}
Học tọa độ tuyệt đối buộc mạng tự ``nhớ'' vị trí hiện tại từ đặc trưng đã chuẩn
hóa: sai số 6 h tăng $\sim$3,5 lần (kém cả persistence) và cần nhiều epoch hơn
(dừng ở epoch 42 thay vì 13--17).

\subsection{Tối ưu huấn luyện}\label{sec:steering}
\textbf{(a) Chỉ dùng thông tin quỹ đạo đã chạm trần.} Trên validation, so với LSTM
gốc (TB 201 km), không thay đổi nào giúp ích: dữ liệu toàn lưu vực (212), lịch sử
48 h (212), loss Haversine (203), mạng lớn/nhỏ hơn (206--207), Transformer (207).
Ensemble trên test chỉ giảm 8 km, KTC 95\% [$-16$, $+2$], không có ý nghĩa.

\textbf{(b) Thêm dòng dẫn đường khí quyển: cải thiện có ý nghĩa.}
\begin{table}[H]\centering\small
\begin{tabular}{lrrrrrr}
\toprule
Test 2021--2024 (3 seed) & 6h & 12h & 24h & 48h & 72h & TB \\
\midrule
LSTM chỉ quỹ đạo (gốc) & 32 & 64 & 131 & 286 & 469 & 197 \\
LSTM + dòng dẫn đường & 30 & 58 & 118 & 270 & 465 & 188 \\
Transformer + dòng dẫn đường & 30 & 57 & 115 & 265 & 450 & 183 \\
\textbf{Ensemble (LSTM + Transformer)} & \textbf{28} & \textbf{54} & \textbf{112} & \textbf{259} & \textbf{443} & \textbf{179} \\
\bottomrule
\end{tabular}
\caption{Hiệu quả của đặc trưng dòng dẫn đường (bão vùng Việt Nam, đủ 48 h lịch sử)}
\end{table}
\begin{itemize}
  \item Ensemble được chọn trên validation; so với mô hình gốc tại 72 h:
        \textbf{$-35$ km, KTC 95\% [$-62$, $-9$]} (bootstrap 44 bão), có ý nghĩa thống kê;
        tại 24 h giảm $\sim$15\%, trung bình các mốc giảm $\sim$9\%.
  \item Gradient boosting + dòng dẫn đường đạt test tốt nhất (TB 176 km) nhưng kém
        hơn trên validation nên \textbf{không được chọn}, giữ đúng quy trình.
\end{itemize}

\subsection{Sai số theo cường độ và nghiên cứu trường hợp}
Áp thấp nhiệt đới khó nhất với persistence (800 km tại 72 h); bão rất mạnh cấp
12--13 có sai số thấp nhất ($\sim$400 km với phương pháp học).

\begin{table}[H]\centering\small
\begin{tabular}{lrrrrr}
\toprule
Bão (thời điểm phát) & k-NN & GBoost & LSTM & Transformer & Persistence \\
\midrule
YAGI 2024 (04/09 06Z) & 419 & \textbf{407} & 541 & 678 & 701 \\
SAOLA 2023 (31/08 12Z) & \textbf{365} & 538 & 534 & 514 & 534 \\
YINXING 2024 (09/11 06Z) & 575 & \textbf{269} & 700 & 550 & 958 \\
CONSON 2021 (10/09 00Z) & 558 & 581 & 625 & 702 & 987 \\
\bottomrule
\end{tabular}
\caption{Sai số 72 h (km) tại thời điểm bão sắp đổ bộ Việt Nam}
\end{table}
Với YAGI, phần lớn mô hình gợi ý bão chuyển hướng lên phía bắc (kiểu quỹ đạo phổ
biến), trong khi YAGI đi thẳng về phía tây; k-NN dựa trên các bão tương tự WUKONG
2000, DAMREY 2005. CONSON có CLIPER-proxy tốt nhất (101 km), một trường hợp đơn lẻ.

\fig{figures/case_WP122024.png}{Nghiên cứu trường hợp YAGI 2024}

\section{Đánh giá hiện trạng}
\textbf{Đã đạt:} pipeline huấn luyện hoàn chỉnh trên dữ liệu thật; chia tập không rò
rỉ; so sánh 7 phương pháp có nhiều seed và kiểm định bootstrap; nón bất định có
độ phủ được kiểm chứng; thí nghiệm ablation và tối ưu huấn luyện có quy trình chọn
mô hình chặt chẽ; mô hình được đóng gói và phục vụ qua API/demo.

\textbf{Nhận xét chính:}
\begin{enumerate}
  \item Học từ dữ liệu giúp rõ ở hạn dài: giảm $\sim$14\% (24 h) đến $\sim$30\% (72 h)
        so với persistence; ở 6 h persistence gần như tối ưu.
  \item Mô hình sâu không vượt trội: với $\sim$7.800 cửa sổ và 4 bước đầu vào, k-NN và
        gradient boosting ngang LSTM/Transformer; mạng nơ-ron overfit nhanh
        (epoch tốt nhất 2--9).
  \item Nút thắt là thông tin, không phải kiến trúc: chỉ khi thêm dòng dẫn đường
        khí quyển mới có cải thiện có ý nghĩa.
  \item k-NN analog có giá trị giải thích, phù hợp tinh thần ``gợi ý'' của đề tài.
\end{enumerate}

\section{Rủi ro và khuyến nghị}
\begin{itemize}
  \item \textbf{Mô hình vận hành chưa dùng dòng dẫn đường:} website/API vẫn phục vụ
        LSTM chỉ quỹ đạo; cần nguồn khí quyển thời gian thực (phân tích GFS).
  \item \textbf{Trường hợp khó:} YAGI 2024 vẫn sai 620--670 km tại 72 h với mọi mô
        hình có dòng dẫn đường; lưới 2,5° chưa đủ nhận diện áp cao cận nhiệt mạnh.
  \item \textbf{Nón bất định đẳng hướng:} sai số dọc/ngang quỹ đạo khác nhau; cần nón
        bất đối xứng.
  \item \textbf{Kiểm định:} kết quả chỉ trên một cách chia theo năm; cần kiểm định chéo
        theo mùa bão.
  \item \textbf{Dữ liệu:} IBTrACS ACTIVE trễ vài ngày; chưa có bản tin thời gian thực.
  \item \textbf{Truyền thông:} giữ disclaimer, nguồn dữ liệu, phiên bản mô hình và
        tham chiếu NCHMF trên mọi giao diện.
\end{itemize}

\section{Kế hoạch ưu tiên tiếp theo}
\begin{enumerate}
  \item Đưa ensemble LSTM + Transformer có dòng dẫn đường vào vận hành, lấy đặc trưng
        từ phân tích GFS tại thời điểm dự báo.
  \item Dùng ERA5 0,25° và đặc trưng dạng ảnh (CNN trên trường gió) để xử lý ca khó
        như YAGI.
  \item Nón bất định bất đối xứng theo dọc/ngang quỹ đạo, hiệu chỉnh theo mốc và cường độ.
  \item Kiểm định chéo theo năm/mùa; chạy \texttt{typhoon-vn monitor} định kỳ để theo
        dõi sai số và độ phủ nón trên bão mới 2025--2026.
  \item Nối registry hiện có (cổng promote, rollback) với MLflow để lưu lịch sử thí
        nghiệm tập trung.
\end{enumerate}

\section{Kết luận}
Dự án đã xây dựng và huấn luyện thành công hệ thống AI gợi ý đường đi của cơn bão
ảnh hưởng Việt Nam trên dữ liệu IBTrACS 1980--2024. Các phương pháp học giảm sai số
72 h khoảng 200 km so với persistence (có ý nghĩa thống kê). Kết quả quan trọng nhất
của phần huấn luyện là: với dữ liệu chỉ có quỹ đạo, LSTM, Transformer, k-NN và
gradient boosting cho độ chính xác tương đương; bổ sung đặc trưng dòng dẫn đường khí
quyển và ensemble giúp giảm thêm 35 km tại 72 h. k-NN analog đóng góp khả năng giải
thích qua danh sách bão tương tự. Hệ thống phù hợp cho học tập và nghiên cứu; trong
sử dụng thực tế, kết quả chỉ là gợi ý bổ sung, không thay thế bản tin của cơ quan
khí tượng có thẩm quyền.

\end{document}
